% command to generate figures
% pdflatex -shell-escape -synctex=1 -interaction=nonstopmode nw2.tex
% 

\documentclass[class=minimal,border=1pt]{standalone}
% \documentclass{article}
\usepackage{tikz}

\usetikzlibrary{calc,matrix,shapes,positioning,fadings}

\tikzfading[name=fade out,
inner color=transparent!0,
outer color=transparent!100]


\newcommand{\convexpath}[2]{
[   
    create hullnodes/.code={
        \global\edef\namelist{#1}
        \foreach [count=\counter] \nodename in \namelist {
            \global\edef\numberofnodes{\counter}
            \node at (\nodename) [draw=none,name=hullnode\counter] {};
        }
        \node at (hullnode\numberofnodes) [name=hullnode0,draw=none] {};
        \pgfmathtruncatemacro\lastnumber{\numberofnodes+1}
        \node at (hullnode1) [name=hullnode\lastnumber,draw=none] {};
    },
    create hullnodes
]
($(hullnode1)!#2!-90:(hullnode0)$)
\foreach [
    evaluate=\currentnode as \previousnode using \currentnode-1,
    evaluate=\currentnode as \nextnode using \currentnode+1
    ] \currentnode in {1,...,\numberofnodes} {
-- ($(hullnode\currentnode)!#2!-90:(hullnode\previousnode)$)
  let \p1 = ($(hullnode\currentnode)!#2!-90:(hullnode\previousnode) - (hullnode\currentnode)$),
    \n1 = {atan2(\y1,\x1)}, 
    \p2 = ($(hullnode\currentnode)!#2!90:(hullnode\nextnode) - (hullnode\currentnode)$),
    \n2 = {atan2(\y2,\x2)},
    \n{delta} = {-Mod(\n1-\n2,360)}
  in 
    {arc [start angle=\n1, delta angle=\n{delta}, radius=#2]}
}
-- cycle
}



% Load the library
\usetikzlibrary{graphs, graphdrawing, external}
\usegdlibrary{circular}
\tikzset{external/system call={lualatex \tikzexternalcheckshellescape --output-format=dvi -halt-on-error -interaction=batchmode -jobname "\image" "\texsource"}}
\tikzexternalize

\begin{document}

% % Figure 0 - Basic Figure with no hidden objects 
% \begin{tikzpicture}[hn/.style={draw,circle,gray, opaque=0.25},
% vn/.style={draw,circle, thick},
% ve/.style={blue,very thick}, 
% he/.style={gray,opaque=0.25,very thick}
% ]
% \graph [
% nodes=vn,
% edges=ve,
% simple necklace layout, 
% ] { 
% 5 [hn], 
% 6 [hn],
% {
% 1 -- 2 , 
% 1 -- 3,
% 4 -- 2, 
% 1 -- 4, 
% 3 -- 4,
% 4 --[he] 5, 
% 3 --[he] 5, 
% 6 --[he] 5, 
% }
% }; 
% \end{tikzpicture}


% First Graph 
\newcommand{\GA}{ 
\begin{tikzpicture}
[
hn/.style={draw,circle,gray},
vn/.style={draw,circle, thick},
ve/.style={blue,very thick}, 
he/.style={gray,opaque=0.25,very thick},
]
\graph [
nodes=vn,
edges=ve,
simple necklace layout, 
] { 

{ \s [hn,source edge style=he,target edge style=he] , 
 %4 [hn,source edge style=he,target edge style=he] , 
 }

{
1 -- 2 , 
1 -- 3,
4 -- 2, 
1 -- 4, 
3 -- 4,
4 --5, 
3 -- 5, 
6 -- 5, 
}
} ;
\end{tikzpicture}
}

% \begin{tikzpicture}[remember picture]
% \node {\includegraphics[height=0.05\lineheight]{rn-network}} ; 
% \end{tikzpicture}


% \begin{tikzpicture}[remember picture]
% \node {\includegraphics[height=0.05\textheight]{rn-network}} ; 
% \fill [blue,path fading=fade out] (-0.5, -1) rectangle (0.5, 1); 
% \end{tikzpicture}


\begin{tikzpicture}[remember picture]
\node {\includegraphics{rn-network}} ; 
\fill [green,path fading=fade out] (-2.5, -0.5) rectangle (-0.5, 1); 
\fill [red,path fading=fade out] (0, -1) rectangle (1.5, 0); 
\end{tikzpicture}



\end{document}
